\subsection{Weil sign, positivity, and all tensor-power sections on the same torus}
\label{subsec:ftb-complex-weil-sections}

Keep exactly the exponential \(e\), lattice \(\Lambda\), frame \(f_i\),
Hermitian form \(H\), unit-modulus function \(\chi\), and integral form
\(E=\operatorname{Im}H|_\Lambda\) constructed in
Proposition~\ref{prop:ftb-actual-appell-humbert-frame}.
Write
\[
 j_\lambda(z)=\chi(\lambda)
       \exp\{\pi H(z,\lambda)+\pi H(\lambda,\lambda)/2\}.
\]
Thus \(f_i(z)=j_\lambda(z)f_i(z+\lambda)\).
The algebraic inputs for this step are precisely
\(\mathrm{Out.WeilSt}\), \srcid{Out.WeilNondegSt}, and
\srcid{Out.GlobalSectionsRankAlgClosed}, in addition to the
multiplication and cube laws already used.
The last law says, for every polarised abelian variety over an
algebraically closed field and every \(n\geq1\),
\(\dim H^0(P,L^n)=n^g\).
We do not replace the source's section comparison by a general
analytic algebraization assertion.

\begin{srlemma}[The actual theta-group scalar and the Weil sign]
\label{lem:ftb-complex-weil-sign}
For \(k\geq1\), \(a,b\in\Lambda\), and
\(x=e(a/k),y=e(b/k)\), the source's algebraic Weil relation is
\[
                 e_{L,k}(x,y)=
                  \exp\{-2\pi i E(a,b)/k\}.                \tag{CW1}
\]
In particular \(E:\Lambda\to\operatorname{Hom}(\Lambda,\mathbb Z)\)
is bijective.
\end{srlemma}

\begin{proof}
Take any algebraic theta-group lift over a test scheme, then restrict
it to a complex point of that test scheme, whose translation is
\(e(a/k)\).  In two charts for \(L\), write \(A_j(z)\) for the value
of its algebraic generator of
\(L_z^k\otimes L_{z+a/k}^{-k}\).  Its transition identity gives
\[
       A_j(z)f_{j_1}(z+a/k)^k
                    =C_a e^{-\pi H(z,a)} f_{j_0}(z)^k      \tag{CW2}
\]
for one scalar \(C_a\) independent of both \(j\) and \(z\).
Here is the source's verification of that scalar.
The quotient
\[
 G_j(z)=
   \frac{A_j(z)f_{j_1}(z+a/k)^k e^{\pi H(z,a)}}
        {f_{j_0}(z)^k}
\]
is holomorphic, and its denominators are nonzero.
The generator and frame transition identities make the \(G_j\)
agree on all overlaps.  Under \(z\mapsto z+\lambda\), their
automorphy factors cancel because
\[
             e^{\pi H(\lambda,a)}=e^{\pi H(a,\lambda)}
\]
for \(\lambda,a\in\Lambda\), by integrality of
\(\operatorname{Im}H(\lambda,a)\).  Hence they glue to a lattice-periodic
entire function and are constant by the fundamental-parallelepiped
and Liouville argument.  Rearranging gives (CW2).

Use (CW2) for the \(a\)-lift at \(z\), the \(b\)-lift at \(z+a/k\),
the \(b\)-lift at \(z\), and the \(a\)-lift at \(z+b/k\).
Cancel the three nonzero intermediate frame values.
The ratio between the two compositions is exactly
\[
 \frac{e^{-\pi H(z,a)}e^{-\pi H(z+a/k,b)}}
      {e^{-\pi H(z,b)}e^{-\pi H(z+b/k,a)}}
       = e^{-2\pi i E(a,b)/k}.                             \tag{CW3}
\]
This fixes the negative sign in (CW1).

This calculation proves the required algebraic generator equality,
not just its value for some chosen lifts.  In the calculation for a
theta-group lift over a test scheme, first restrict that lift to a
complex point of the test scheme; (CW2)--(CW3) are then calculations
on the actual complex variety \(P\).
For the predicate \(\mathrm{WeilIs}\), the source applies
\(\mathrm{WeilSt}\) over the base \(S=\operatorname{Spec}\mathbb C\).
It obtains its lifts and its Weil relation after restriction to an
open \(V\subset S\).  The space on which the resulting regular
generator equality is checked is therefore \(V\times_{\mathbb C}P\).
This space is smooth, reduced and locally of finite type over
\(\mathbb C\).  Formula (CW3), at all its complex points, determines
the scalar in that relation.  Equality of the regular functions
follows on each affine open: their difference lies in every maximal
ideal of a finitely generated complex algebra, hence in its zero
nilradical.  No reducedness assertion about a product with an
arbitrary test scheme is used.
Restriction, pullback, multiplication of generators and constant
units all have exactly their evaluated products in this argument.
This is the source predicate \(\mathrm{WeilIs}\) with scalar (CW1).

To prove perfection, suppose \(k\mid E(a,b)\) for every \(b\in\Lambda\).
Every \(k\)-torsion point is \(e(b/k)\): lift it to \(w\in V\), so
\(kw\in\Lambda\), and take \(b=kw\).
Formula (CW1) therefore says that \(e(a/k)\) pairs trivially with
all \(k\)-torsion points.  The supplied algebraic nondegeneracy
\(\mathrm{WeilNondegSt}\) gives \(e(a/k)=0\).
Thus \(a/k\in\Lambda\), proving
\[
 \bigl(k\mid E(a,b)\text{ for every }b\bigr)
                            \Longrightarrow a\in k\Lambda. \tag{CW4}
\]

For clarity, the source's integral algebra from (CW4) is as follows.
The lattice is finite free over \(\mathbb Z\).
If \(E(a,-)=0\), (CW4) makes every coordinate of \(a\) divisible
by every positive \(k\), so \(a=0\).  Thus \(E\) is injective.
Its cokernel has rank zero, since source and target have equal
finite rank.  For any functional \(\ell\), choose a nonzero integer
\(c\) and \(x\in\Lambda\) with \(E(x,-)=c\ell\).
Then \(E(cx,-)=c^2\ell\), so (CW4) gives \(cx=c^2a'\) for some
\(a'\).  Applying \(E\) and cancelling \(c^2\ne0\) gives
\(E(a',-)=\ell\).  This proves surjectivity.
\end{proof}

For \(n\geq0\), let \(\mathcal T_n\) be the vector space of entire
functions satisfying
\[
                    \theta(z+\lambda)=j_\lambda(z)^n\theta(z).
                                                                    \tag{CW5}
\]
The actual section-to-function map is
\[
            \theta_s(z)=s_i(e(z))/f_i(z)^n
                         \quad\text{whenever }e(z)\in U_i. \tag{CW6}
\]
It is independent of \(i\), by the matching section and frame
transition identities.  It is entire because its local denominator
never vanishes, and (CW5) follows by evaluating at \(z+\lambda\).
The same nonvanishing proves uniqueness of \(\theta_s\).
Conversely, two algebraic sections with the same function have the
same values at every complex point, by surjectivity of \(e\).
The reduced finite-type argument in the proof of
Lemma~\ref{lem:ftb-complex-weil-sign}, applied to each section chart,
therefore makes the two sections equal.  Thus (CW6) is an injective
linear map \(H^0(P,L^n)\to\mathcal T_n\) for every \(n\).

\begin{srlemma}[Positivity from the source's nonzero theta function]
\label{lem:ftb-complex-theta-positivity}
The Hermitian form \(H\) is positive definite.
\end{srlemma}

\begin{proof}
The rank law for \(n=1\) supplies a nonzero algebraic section.
Its function \(\theta\) under (CW6) is nonzero by injectivity.
The continuous function
\[
           |\theta(z)|\exp\{-\pi H(z,z)/2\}
\]
is lattice-periodic: \(|\chi|=1\), and the real part of the
Hermitian automorphy exponent gives the precise quadratic
difference.  It is bounded on a closed fundamental parallelepiped.
Consequently, for some real \(C\geq0\),
\[
                  |\theta(z)|\leq C e^{\pi H(z,z)/2}.       \tag{CW7}
\]
Here \(H(z,z)\) is real.

If \(H(w,w)<0\), then for each fixed \(z\) the entire one-variable
function \(t\mapsto\theta(z+tw)\) has, by (CW7), a bound
\[
 C\exp\{\tfrac\pi2(A+2B|t|-\delta|t|^2)\},
 \qquad \delta=-H(w,w)>0,\quad B=|H(z,w)|.
\]
It tends to zero at infinity.  Liouville's theorem makes it identically
zero.  Varying \(z\) would give \(\theta=0\), a contradiction.
Hence \(H\) is positive semidefinite.

The source next removes its possible nullspace using \(E\).
If \(H(z,z)=0\), positivity of
\(H(z+isv,z+isv)\) for every real \(s\) gives
\[
                 0\leq 2s\,\operatorname{Im}H(z,v)
                                  +s^2H(v,v).
\]
Choose \(s=-c/(d+1)\), where
\(c=\operatorname{Im}H(z,v)\) and \(d=H(v,v)\geq0\).
The resulting expression is
\(-c^2(d+2)/(d+1)^2\), so \(c=0\).
Choose an integral basis of \(\Lambda\), which is a real basis of
\(V\).  In these coordinates the last equality for each lattice
basis vector says that the coordinate row of \(z\) lies in the
left kernel of the matrix of \(E\).  That matrix has nonzero
determinant by Lemma~\ref{lem:ftb-complex-weil-sign}; it remains
invertible over \(\mathbb R\).  Therefore \(z=0\).
This proves positive definiteness.
\end{proof}

\begin{srlemma}[The source's finite Fourier detector]
\label{lem:ftb-complex-theta-fourier-detector}
For \(n\geq1\), every linearly independent family in
\(\mathcal T_n\) has at most \(n^g\) members.
\end{srlemma}

\begin{proof}
We retain the actual symplectic normal form and its gauge.
The full lattice has rank \(2g\).  The alternating perfect form
\(-E\) has an integral symplectic basis \((a_i,b_i)\), so
\(E(a_i,b_j)=-\delta_{ij}\).
The source obtains this basis by the following integral elimination.
Among ideals generated by a value of an alternating form on the
current submodule, choose a maximal one.  A Bezout combination
shows that its generator \(d=B(e,f)\) divides \(B(e,y)\),
\(B(x,f)\), and every value on the common orthogonal complement.
Subtract the corresponding multiples of \(e,f\) to decompose each
vector into their span and that complement.  Induction on rank
gives orthogonal pairs spanning the module.  Their nonzero \(d_i\)
give linear independence; surjectivity of the form forces each
\(d_i\) to be a unit by applying it to a basis coordinate functional.
Rescale the second member of each pair by this unit.  The resulting
basis is symplectic.  This is applied to \(-E\), explaining the
sign of \(E(a_i,b_j)\).

Let \(A\) have columns \(a_i\).  The isotropic relation makes
\(A^tH\overline A=Y\) a real matrix.  It is invertible:
otherwise a real kernel vector of \(Y\), positivity of \(H\),
and the pairing with every \(b_j\) would force that vector to be
zero.  Thus \(A\) is a complex isomorphism.  Write the matrix of
the second half of the basis as \(A\tau\).
The pairing equations give
\[
             Y>0,\qquad Y\,\operatorname{Im}\tau=1,\qquad
             \tau^t=\tau,\qquad \operatorname{Im}\tau>0 .
\]
The symmetry follows by substituting the zero pairings between
the \(b_i\)'s, separately in real and imaginary parts.

For appropriate complex linear \(c\) put
\(Q(z)=\pi z^tYz/2+c(z)\).
Choose \(c(e_i)=\log\chi(a_i)\).  Direct subtraction of \(Q\)
under \(z\mapsto z+e_i\) and \(z\mapsto z+\tau e_j\) shows that
\[
                   F(z)=\theta(Az)e^{-nQ(z)}
\]
is entire, \(1\)-periodic in each coordinate, and
\[
          F(z+\tau e_j)=\beta_j^n
                    e^{-2\pi i n z_j}F(z),\qquad \beta_j\ne0. \tag{CW8}
\]
The map \(\theta\mapsto F\) is injective since \(A\) is invertible
and the gauge exponential never vanishes.

For such an \(F\), define the actual unit-cube Fourier coefficients
\[
             c_m(F)=\int_{[0,1]^g}F(x)e^{-2\pi i m\cdot x}\,dx,
                      \qquad m\in\mathbb Z^g .
\]
The source's shift identity is
\[
 \int_{[0,1]^g}F(x+w)e^{-2\pi i m\cdot x}\,dx
                      =e^{2\pi i m\cdot w}c_m(F).           \tag{CW9}
\]
Indeed the integral of an entire periodic function over the real
unit cube is unchanged by a real translation (Haar translation
on the unit torus).  Differentiation under the integral on this
compact cube makes the integral an entire function of the complex
translation.  Real translation invariance then makes it constant.
Apply this to \(F(z)e^{-2\pi i m\cdot z}\) to get (CW9).

Combining (CW8)--(CW9) gives
\[
       \beta_j^n c_{m+ne_j}(F)
                         =e^{2\pi i m\cdot\tau e_j}c_m(F).  \tag{CW10}
\]
Both multipliers are nonzero.  Vanishing of the coefficients
with \(0\leq m_j<n\) therefore implies vanishing for every
\(m\in\mathbb Z^g\): use integer division of each coordinate
by \(n\) and move by the steps \(\pm ne_j\).
If all coefficients vanish, \(F=0\).  To see this last assertion
on any complex translate of the real torus, (CW9) says that all
its Fourier coefficients vanish.  Fourier uniqueness for continuous
functions on the compact real torus gives zero on that translate,
in particular at its origin.  The translate was arbitrary.

Thus the \(n^g\) coefficients in the residue box define an injective
linear map on this function space into \(\mathbb C^{n^g}\).
Applying this map to a linearly independent family proves the claim.
\end{proof}

\begin{srproposition}[The full fixed-fibre realization]
\label{prop:ftb-full-fixed-fibre-theta-realization}
With the supplied multiplication, cube, Weil, nondegeneracy and
algebraic section-rank laws, the single tuple
\((e,\Lambda,f_i,H,\chi,E)\) above satisfies the full
\(\mathrm{ThetaUnifAV}\) conclusion.  In particular \(H>0\),
\(E\) is perfect, (CW1) has its displayed negative sign, and for
every \(n\geq0\) the rule (CW6) is a bijection between the
algebraic sections of \(L^n\) and \(\mathcal T_n\).
The tangent comparison, linear endomorphism lifts, original
rigidification, and exponential are the ones previously constructed.
\end{srproposition}

\begin{proof}
For \(n\geq1\), the algebraic rank law supplies \(n^g\) independent
sections.  Their functions are independent: a zero linear
combination of functions comes from the same linear combination
of sections by (CW6), so injectivity gives a zero combination of
sections.  Lemma~\ref{lem:ftb-complex-theta-fourier-detector}
then makes these functions span \(\mathcal T_n\).
Otherwise adjoining one function outside their span would give
\(n^g+1\) independent functions.  The corresponding linear combination
of the original algebraic sections gives the required section;
injectivity proves uniqueness.

For \(n=0\), (CW5) says simply that \(\theta\) is lattice-periodic.
It is bounded on a fundamental parallelepiped and hence constant
by Liouville.  The algebraic constant section of \(L^0=\mathcal O_P\)
has exactly this function.  This also gives uniqueness.
The remaining conclusions are precisely the preceding lemmas
and the exponential and frame constructions, whose exponentials
and lattices were unchanged.
\end{proof}

\paragraph{Exact frozen-source ledger.}
Unless otherwise indicated the following paths are under
\srcid{src/Development/AbelianVarietyTheory/Proofs/ToricGeometry/},
at tag \srcid{alpoge-source-7cbf7a0c}.
The scalar comparison (CW2) is
\srcid{ExistsConstOfCover.lean}
(\srcid{352bc50a120c49044d56f2cb6d27897e02e9ded2}),
declaration \srcid{ThetaGroupFrame.thetaGrp\_in\_frame}.
The four-frame cancellation and algebraic Weil relation are
\srcid{WeilIsExpRiemannForm.lean}
(\srcid{1eb1a1affef7a79c6e3ba7d22535430afa9f5930}).
The divisibility step from nondegeneracy is
\srcid{RiemannFormBijectiveOfWeilNondeg.lean}
(\srcid{1f530935aa3e1561f47e7164efbc1e450969a70a});
the finite-free rank and \(c^2\) argument is
\srcid{BilinFormBijectiveOfNondegenerateMod.lean}
(\srcid{996fa526419927f2aa701e51691e6593347aa79e}).
Equality of regular functions from their complex values uses
\srcid{EqZeroOfForallAlgHomComplex.lean}
(\srcid{292d8cb01972525a6c0040b1449efb5090dcca4f}) and
\srcid{IsReducedSchemeOfSmoothOfIsIntegrallyClosed.lean}
(\srcid{14b7a75b14c084dda1e8a2b2b06dd4d26a0c3fba}).

Construction and injectivity of (CW6) are
\srcid{SectionsPowThetaFunctions.lean}
(\srcid{54dfefb784a86c4772781b543aa65b04d961f540}).
The nonzero rank-one section and positivity assembly are
\srcid{ExistsIndepSections.lean}
(\srcid{fa66eca4089a8ad5a3bf75a329926da7b4bc58d8}) and
\srcid{IsAHOfPolarisation.lean}
(\srcid{22ea6df6cd8176d9647db8c52e101b63203fd5a2});
the periodic metric bound, negative-direction argument and
nullspace removal are
\srcid{ReFormPosOfThetaNeZero.lean}
(\srcid{bda58e9df3981e4de3d25e29bb8531c5f4f93e1e}).

The integral symplectic elimination is
\srcid{../CovolumeBounds/DualRestrictSurjectiveOfIsTorsionFreeQuotient.lean}
(\srcid{369a733eead2b44d17430d13ba108b43be4e8f81}).
Its application, the exact \(-E\) sign, period matrix and gauge are
\srcid{IsUnitOfLagrangian.lean}
(\srcid{b2457d0adf49b8651e52f5bf0a955de91cbdefda}),
declaration \srcid{SiegelNormalForm.exists\_normal\_form}.
The complex Fourier shift is
\srcid{FourierCoeffUnitCubeTranslateOfEntirePeriodic.lean}
(\srcid{90e515c3c1229348b46c587391faf4f2f1dfbedd});
Fourier uniqueness on translates is
\srcid{EntirePeriodicEqZeroOfFourierCoeffEqZero.lean}
(\srcid{ba105521538d00b051b34f19676d77e0f374ca90}).
The residue-box argument is
\srcid{CardLePowOfLinearIndependent.lean}
(\srcid{20085cbe9a9a42f684df9924f67d1e8fcc10a49f});
the normal-form transport is
\srcid{CardLePowOfLinearIndependentThetaPrincipal.lean}
(\srcid{fe36db1b4f2c9d4f78040330571ea9d7ca87ef58}).
Surjectivity including \(n=0\) is
\srcid{MemSpanOfCardLe.lean}
(\srcid{f21dc5ff7a5efc6072c172914c28130acaafe764}).
Final same-tuple assembly is
\srcid{ThetaUnifAVOfParts.lean}
(\srcid{d260f9fb5b92e8635a1893372bb0ad47bdacb5c6}),
called in the fixed order by
\srcid{ThetaUnifAVOfLaws.lean}
(\srcid{3e496ea337ef9e519bbf43dfe91ef52009f6a960}).
This completes that producer from its five displayed algebraic laws.
The construction of families of period matrices and the reverse
construction of an algebraic variety from a given period matrix
are separate later producers.
